% SPDX-License-Identifier: Apache-2.0
% covers: LinePair
\datasheet{LinePair}{The pixel-clock end of a scanout from memory: two
lines, one shown while the other fills, each asked for one or two rows
ahead.}

\dsfacts{\code{lib/parts/src/scanout.rs}}%
{\code{txhdl\_parts::scanout::LinePair<LEN, AW, ROWS, TOTAL, STRIDE, LO,
HI, C>}}%
{\code{//docs:examples}}

\subsection*{Function}

A video peripheral that shows a frame from memory cannot wait for
memory pixel by pixel: the beam does not stop. \code{LinePair} sits on
the pixel clock beside the raster and holds two lines. The beam reads
one while the other fills, and they change places at the start of each
line. At that start it asks for the line after the one about to be
shown, so each line has the entire duration of the preceding line to
arrive. With \code{two} high at the vertical sync it asks two rows
ahead for the frame that follows, so a line has two rows to come, and
its words wait in the crossing until the row before its own
(issue 1523).
It asks for nothing until \code{show} is high, and then it asks for a
frame's first line, from the base, before any other. Out of a reset
the next line's address is zero, and the line there is the boot
memory's, whose one-beat answer to a burst hung the fetch for good on
the board (issue 1178). So, one row ahead, the picture starts with
the first frame whose last row begins after \code{show} rises.
The bus side of the same scanout is \code{ScanFetch}, and between the
two the request and the words cross clocks through \code{ChanCdc}
(issue 151).

The frame's base address is a register taken at the vertical sync, so
a host that draws into one buffer and shows another flips them with
one write. A column the beam shows before its word has arrived sets
\code{starved}, a sticky bit that \code{clear} resets, so a run can
show that no line starved. The pair keeps the addresses of the lines
it has asked for and not had whole, at most four, with the row each is
for and the time it was asked at. A line that gets no word for two
line times in a row sets \code{stuck}, sticky too, and puts its
address on \code{stuck\_at}, so a fetch that hangs says so rather than
leaving a black screen and a clear \code{starved} (issue 1197).
\code{worst} is the longest a line took from being asked for to its
last word, in pixels, and \code{lates} counts the lines not whole when
their rows began; \code{clear} resets both (issue 1523).

The frame's memory is the bytes from \code{LO} up to but not including
\code{HI}. A line not wholly inside them is not asked for at all: it
sets \code{stuck} with its address and shows \code{LATE}, and the
frame goes on (issue 1382). A base gone wrong once sent a line's burst
into the serial port's page, and the core's prints stopped behind it.

\subsection*{Parameters}

\begin{center}\footnotesize
\begin{tabular}{@{}lp{0.68\textwidth}@{}}
\toprule
Parameter & Meaning \\
\midrule
\code{LEN} & The words in a visible line, at most \code{1 << AW}. \\
\code{AW} & The width of a column. \\
\code{ROWS} & The visible rows of a frame. \\
\code{TOTAL} & The rows of a frame, blanking included. \\
\code{STRIDE} & The bytes from one line's start in memory to the next's. \\
\code{LO}, \code{HI} & The frame's memory: the bytes from \code{LO} up
to but not including \code{HI}. \\
\code{C} & The pixel clock. \\
\bottomrule
\end{tabular}
\end{center}

The tables below use \code{LinePair<8, 3, 4, 6, 32, 0, 1024, ClkPix>},
the pair of \code{ex\_scanpair}: lines of eight words, four visible
rows of six, and a memory of 1024 bytes from zero. The flagship's,
inside \code{ScanVideo}, is \code{LinePair<640, 10, 480, 525, 4096,
0x4000\_0000, 0x8000\_0000, \ldots>}: the board's gigabyte of DDR3.

\dstables{LinePair}

\subsection*{Behaviour}

\begin{itemize}
\item A word is taken whenever one is offered, into the line filling
at the next column. The pair is what the fetch is backpressured by,
and a line longer than \code{LEN} loses its tail rather than stalling.
\item At a line's start (\code{line} high) the two lines change
places, and the line filling's count starts again. A word taken in
that same cycle counts in the line just filled.
\item At a line's start a request goes out on \code{req}, if it has
room: the address of the line for the row one ahead, or two ahead
when \code{two\_ahead} is set. When that row is the frame's first, it
is the frame's first line, at the base taken at the last vertical
sync, or at \code{base} itself when the sync is that same edge.
\item \code{two\_ahead} is \code{two} as it was at the last vertical
sync. Two rows ahead, a word is not taken while the oldest line owed
is for a row other than the one the half filling is for, unless that
line is late or behind: its words wait in the crossing until the row
before its own.
\item A line not wholly inside \code{LO} to \code{HI} is not asked
for, and is not owed. If \code{stuck} is clear it sets it, with that
line's address on \code{stuck\_at}. The next line's address moves on
as if it had been asked for, so the rows after it are asked where they
belong, and its row shows \code{LATE}.
\item While \code{show} is low nothing is asked for. The frame's first
line is asked for at the start of the row before the frame's first, or
two rows before, with \code{show} high, and the lines after it only
once it has been, so the first request after a reset or a hide is
always the base's line.
With \code{show} low the pair disarms: it asks for nothing more, and
does not set \code{starved}. Shown again, it starts from the base's
line as out of a reset.
\item \code{pix} is the word at \code{col} in the line not filling, a
cycle after \code{col} names it, or \code{LATE}, magenta
(\code{0x00ff00ff}), if that word has not arrived (issue 1209). No
picture is drawn in magenta, so a late column shows as itself rather
than as the word the line before left there.
\item A word belongs to the line it was asked for, which the lines owed
say (issue 1233). A line that starts before it has come whole is
\code{late}, and the rest of its words go into the half the beam
reads, so its columns fill in as they come. Lines owed whose time has
passed are counted in \code{behind}, and the rest of their words are
dropped. Neither lands in the half filling, where the next row would
show it as its own.
\item Once the frame's first line has been asked for, a visible column
at or past the words that line had when it was swapped in sets
\code{starved}, until \code{clear}.
\item A line asked for goes after the lines still owed, with its row
in \code{g0} to \code{g3} and the \code{tick} it was asked at in
\code{t0} to \code{t3}, and no line is asked for while four are owed.
A word counts against the oldest owed, and that line's last word takes
it off.
\item \code{tick} counts the edges of \code{C}. When a line's last word
comes, the ticks since it was asked for replace \code{worst} if they
are more. At the start of each visible row, once armed, a line not
whole adds one to \code{lates}, which stops at its largest.
\code{clear} sets both to zero.
\item Two line starts in a row with lines owed and no word come set
\code{stuck}, and \code{stuck\_at} is the oldest owed line's address
then; while \code{stuck} is set it does not change. \code{clear}
clears \code{stuck} and leaves \code{stuck\_at}. The pair goes on
asking.
\end{itemize}

\subsection*{Verification}

\begin{itemize}
\item \code{ex\_scanpair} runs \code{LinePair} on a pixel clock
unrelated to the bus clock, with \code{ScanFetch}, \code{LineFetch},
\code{AxiHost}, \code{AxiPer}, a memory and \code{ChanCdc} both ways.
The memory stalls its beats, standing in for a refresh. Three frames
with a short stall in every line show every pixel as the word at its
place and never set \code{starved}; a fourth with one stall longer than
a line sets it, and it stays set until it is cleared.
\item \code{//docs:scanpair\_sim\_linepair\_linepair\_tb\_test} and
\code{//docs:scanpair\_vsim\_linepair\_test} simulate the netlist
against that run under nvc and Verilator.
\item In \code{//lib/parts:parts\_test}, the flagship's pair and
raster:
\begin{itemize}
\item \code{a\_frame\_with\_its\_words\_on\_time\_has\_no\_late\_line}
gives the words as fast as the pair takes them, one row ahead and two,
and from the second frame sees no late line, no starved column and no
\code{stuck}.
\item \code{the\_blanking\_columns\_read\_nothing\_past\_the\_line}
runs two lines of 640 words in 800 columns (issue 1194).
\end{itemize}
\item \code{//cpu/vreteno:board\_test} runs the flagship's line,
\code{LinePair<640, 10, 4, 6, 4096, LO, HI, Pix>}, on a pixel clock a
quarter of the board's, against the board:
\begin{itemize}
\item \code{a\_base\_outside\_the\_memory\_is\_never\_fetched} gives it
the base \code{0x3000}, the serial port's page: no line is asked for,
\code{stuck} is set at \code{0x3000}, and the core's greeting comes
out whole.
\item \code{a\_stalled\_read\_is\_late\_one\_row\_ahead\_and\_not\_two}
stalls one line's read for 4000 cycles: one row ahead, four lines are
late and two rows show \code{LATE}; two rows ahead, one line is late
and no column starves. Either way every row shows its own line.
\item \code{a\_flipped\_frame\_is\_shown\_one\_row\_ahead\_and\_two}
flips between two bases in the middle of every frame, and one
row ahead and two every row shows its own line from the frame the pair
took last.
\end{itemize}
\end{itemize}

\subsection*{Used by}

The example \code{ex\_scanpair}, and \code{ScanVideo}, which the
flagship holds in its video slot (issue 151).

\subsection*{Limits}

One line a request, and a line is \code{LEN} words; there is no
horizontal offset or scaling. A column of a line that arrives late
shows \code{LATE} until its word comes, and \code{starved} says that
it happened; a line that never comes sets \code{stuck}. A line outside
\code{LO} to \code{HI} is never shown. It runs on one clock,
\code{C}.
